% !TeX root = ../机器学习笔记.tex
% 由主文档通过 \input 载入；请编译主文档。
\setcounter{chapter}{2} % 保留原书章号。
\chapter{现有分章材料与全书的对应关系}
\label{app:ml-materials}

\section{已有材料索引}

分章材料现位于本课程的“笔记与习题／教材研读／机器学习分章研读”目录中，各章单独存放，目录名包含章号与主题。下表按当前主题结构给出阅读对应；对应关系表示内容关联，不表示这些独立材料已经逐式核验或被直接合并进本书。
\begin{longtable}{@{}L{2.7cm}L{5.3cm}L{5.7cm}@{}}
  \toprule
  分章材料 & 现有主题 & 本书对应位置 \\
  \midrule
  \endhead
  第 1 章 & 统计学习与监督学习概论 & 第\ref{chap:ml-tasks}、\ref{chap:ml-evaluation}、\ref{chap:ml-generalization}章。\\
  第 2 章 & 感知机 & 第\ref{chap:ml-logistic}章。\\
  第 3 章 & $k$ 近邻法 & 第\ref{chap:ml-neighbors}章。\\
  第 4 章 & 朴素贝叶斯法 & 第\ref{chap:ml-generative}章。\\
  第 5 章 & 决策树 & 第\ref{chap:ml-trees}章。\\
  第 6 章 & 逻辑回归与最大熵模型 & 第\ref{chap:ml-logistic}章。\\
  第 7 章 & 支持向量机 & 第\ref{chap:ml-kernels}章。\\
  第 8 章 & 提升方法 & 第\ref{chap:ml-ensembles}章。\\
  第 9 章 & EM 算法及其推广 & 第\ref{chap:ml-em}章。\\
  第 10 章 & 隐马尔可夫模型 & 第\ref{chap:ml-graphical}章。\\
  第 11 章 & 条件随机场 & 第\ref{chap:ml-graphical}章。\\
  第 12 章 & 监督学习方法总结 & 第二篇与附录\ref{app:ml-methods}。\\
  第 13 章 & 无监督学习概论 & 第三篇与第\ref{chap:ml-tasks}章。\\
  第 14 章 & 聚类方法 & 第\ref{chap:ml-clustering}章。\\
  第 15 章 & 奇异值分解 & 第\ref{chap:ml-linear-algebra}、\ref{chap:ml-dimension}章。\\
  第 16 章 & 主成分分析 & 第\ref{chap:ml-dimension}章。\\
  第 17 章 & 潜在语义分析 & 第\ref{chap:ml-text}章。\\
  第 18 章 & 概率潜在语义分析 & 第\ref{chap:ml-text}章。\\
  第 19 章 & 马尔可夫链蒙特卡罗法 & 第\ref{chap:ml-bayesian}章。\\
  第 20 章 & 潜在狄利克雷分配 & 第\ref{chap:ml-text}章。\\
  第 21 章 & PageRank 算法 & 第\ref{chap:ml-graphs}章。\\
  第 22 章 & 无监督学习方法总结 & 第三篇、第五篇与附录\ref{app:ml-methods}。\\
  \bottomrule
\end{longtable}

\section{结合分章材料复习}

复习时先用本书的任务、模型、目标和算法框架定位方法，再查阅对应分章材料中的长推导；不同材料可能采用不同标签编码、损失平均和正则化系数，转用公式前须统一这些约定。对每种方法，至少独立完成一遍核心推导和一个数值例题，再用独立验证比较简单基线，避免将熟悉公式形式误认为已经理解假设。

学习理论的进一步证明可查阅文献\cite{ml-uml}，凸优化与对偶可查阅文献\cite{ml-convex}，概率模型与推断可查阅文献\cite{ml-prml}。统计建模方法的系统比较见文献\cite{ml-esl}；有限链与混合性质见文献\cite{ml-markov}，表格 Q-learning 的原始收敛论证见文献\cite{ml-qlearning}。这些来源用于补充本书明确标注的证明范围，不表示本书复现了其中全部理论和实验。




